% =============================================================================
% Worked solutions for Listas/lista-r2-valuation.md (Ricca R2 -- Valuation).
% Statements are not restated: each section cites the exercise by number and page.
% Compile: pdflatex (twice); check: pdffonts (Type 1 only, uni = yes).
% Numbers: python -I Listas/lista-r2-valuation-solutions_check.py  -> PASS
% =============================================================================
\documentclass[11pt,a4paper]{article}
\usepackage[utf8]{inputenc}
\usepackage[T1]{fontenc}
\usepackage{lmodern}   % Type 1 vector fonts (replaces PK bitmaps)
\usepackage{cmap}      % ToUnicode CMap -> copyable text and accents
\usepackage[brazilian,english]{babel}
\usepackage{amsmath,amssymb,amsfonts,amsthm}
\usepackage[a4paper,margin=2.2cm]{geometry}
\usepackage{booktabs}
\usepackage{enumitem}
\usepackage{fancyhdr}
\usepackage{xcolor}
\usepackage{microtype}

\pagestyle{fancy}\fancyhf{}
\lhead{\footnotesize Finan\c{c}as II -- Corporativas (MPE Insper)}
\rhead{\footnotesize List R2 --- Worked solutions}
\cfoot{\footnotesize\thepage}
\renewcommand{\headrulewidth}{0.3pt}
\setlength{\headheight}{14pt}

\makeatletter
\renewcommand{\@seccntformat}[1]{\csname the#1\endcsname.\quad}
\makeatother
\numberwithin{equation}{section}
\setlist{topsep=2pt,itemsep=1pt}

\newcommand{\step}[2]{\par\medskip\noindent\textbf{Step #1 --- #2.}~}
\newenvironment{notation}{\par\noindent\textbf{Notation.}\begin{itemize}[leftmargin=1.4em]}{\end{itemize}}
\newcommand{\intuition}[1]{\par\medskip\noindent\textbf{Intuition.}~#1}
\newcommand{\refline}[1]{\par\smallskip\noindent{\small\color{black!70}\textbf{Ref:} #1}}
\newcommand{\bookcheck}[1]{\par\smallskip\noindent\textbf{Book answer.}~#1}
\newcommand{\remark}[1]{\par\smallskip\noindent\textit{Remark.}~#1}
\newcommand{\pt}{\foreignlanguage{brazilian}}

\begin{document}

\begin{center}
  {\Large\bfseries List R2 --- Valuation: free cash flows, discount rates, and the NPV rule}\\[4pt]
  {\large Worked solutions}
\end{center}

\noindent Statements: \texttt{Listas/lista-r2-valuation.md} (copied verbatim from the books; not
restated here). Each section cites its exercise by book, number and \emph{printed} page of the
local editions: Berk, DeMarzo \& Harford, \emph{Fundamentals of Corporate Finance}, 2nd ed.\
(2011); Koski, Higgins \& Mitton, \emph{Analysis for Financial Management}, 13th ed.\ (2023);
Grinblatt \& Titman, \emph{Financial Markets and Corporate Strategy}, 2nd ed.\ (2002); Ross,
Westerfield, Jordan \& Lamb, \emph{Fundamentos de Administra\c{c}\~ao Financeira}, 9th ed.\ (2013);
Amaro de Matos, \emph{Theoretical Foundations of Corporate Finance} (2001). Rates are written as
decimals inside equations. Every number below is reproduced by
\texttt{Listas/lista-r2-valuation-solutions\_check.py}.

% =============================================================================
\section{NPV of a two-year investment --- Berk 8.2, p.~241}

\begin{notation}
  \item $C_0=100{,}000$: investment paid today, in dollars ($t=0$).
  \item $C_1=80{,}000$, $C_2=30{,}000$: cash inflows at $t=1$ and $t=2$, in dollars.
  \item $r=0.09$: cost of capital per year, $r>0$.
  \item $\mathrm{NPV}$: net present value at $t=0$, in dollars, any real number.
\end{notation}

\step{1}{NPV definition} The NPV is the present value of all cash flows, outflows negative:
\begin{equation}\label{eq:1npv}
  \mathrm{NPV} = -C_0 + \frac{C_1}{1+r} + \frac{C_2}{(1+r)^2}.
\end{equation}

\step{2}{Substitute the data} Placing the numbers in \eqref{eq:1npv}:
\begin{align}
  \mathrm{NPV} &= -(100{,}000) + \frac{(80{,}000)}{(1.09)} + \frac{(30{,}000)}{(1.09)^2} \notag\\
               &= -100{,}000 + 73{,}394.50 + 25{,}250.40 \notag\\
               &= -1{,}355.10. \label{eq:1res}
\end{align}

\step{3}{Decision} By \eqref{eq:1res} the NPV is negative: the inflows are worth \$98{,}644.90
today, less than the \$100{,}000 they cost. Reject the investment.

\intuition{The inflows sum to \$110{,}000, a 10\% total return over two years, which is less than
9\% per year compounded. Money returned later is worth less, so the project destroys
\$1{,}355.10 of value relative to investing at the cost of capital.}

\refline{Berk, \S8.1--8.2 (The NPV decision rule; using the NPV rule), pp.~212--216.}

% =============================================================================
\section{CAPM: Intel and Boeing --- Berk 12.24, p.~377}

\begin{notation}
  \item $r_f=0.04$: risk-free rate; $E[R_m]=0.10$: expected market return; $E[R_m]-r_f$: market risk premium.
  \item $\beta_I=1.6$, $\beta_B=1.0$: betas of Intel and Boeing stock, real numbers.
  \item $E[R_I]$, $E[R_B]$, $E[R_P]$: expected returns of Intel, Boeing and the portfolio $P$.
  \item $w_I=0.6$, $w_B=0.4$: portfolio weights, $w_I+w_B=1$; $\beta_P$: beta of $P$.
\end{notation}

\step{1}{The CAPM (security market line)} For any security $i$,
\begin{equation}\label{eq:2sml}
  E[R_i] = r_f + \beta_i\,(E[R_m]-r_f).
\end{equation}

\step{2}{(a) Intel} Substituting $\beta_I$ into \eqref{eq:2sml}:
\begin{align}
  E[R_I] &= (0.04) + (1.6)(0.10-0.04) \notag\\
         &= 0.04 + 0.096 \notag\\
         &= 0.136 = 13.6\%. \label{eq:2I}
\end{align}

\step{3}{(b) Boeing} Substituting $\beta_B$ into \eqref{eq:2sml}:
\begin{align}
  E[R_B] &= (0.04) + (1.0)(0.10-0.04) \notag\\
         &= 0.04 + 0.06 \notag\\
         &= 0.10 = 10\%. \label{eq:2B}
\end{align}

\step{4}{(c) Portfolio beta} Beta is linear in portfolio weights:
\begin{align}
  \beta_P &= w_I\beta_I + w_B\beta_B \label{eq:2bp}\\
          &= (0.6)(1.6) + (0.4)(1.0) \notag\\
          &= 0.96 + 0.40 \notag\\
          &= 1.36. \label{eq:2bpres}
\end{align}

\step{5}{(d) Portfolio expected return, two ways} First, substituting $\beta_P$ from
\eqref{eq:2bpres} into \eqref{eq:2sml}:
\begin{align}
  E[R_P] &= (0.04) + (1.36)(0.06) \notag\\
         &= 0.04 + 0.0816 \notag\\
         &= 0.1216 = 12.16\%. \label{eq:2p1}
\end{align}
Second, as the weighted average of \eqref{eq:2I} and \eqref{eq:2B}:
\begin{align}
  E[R_P] &= w_I E[R_I] + w_B E[R_B] \notag\\
         &= (0.6)(0.136) + (0.4)(0.10) \notag\\
         &= 0.0816 + 0.04 \notag\\
         &= 0.1216 = 12.16\%. \label{eq:2p2}
\end{align}
The two routes agree because both \eqref{eq:2sml} and \eqref{eq:2bp} are linear in the weights.

\intuition{Under the CAPM only market risk is priced, and it is measured by beta. Intel moves 1.6
times as much as the market, so it earns 1.6 risk premiums. A portfolio is priced by its
average beta, so averaging returns or averaging betas gives the same answer.}

\refline{Berk, \S12.1 (expected return of a portfolio), p.~347; \S12.4 (the CAPM), pp.~366--372.}

% =============================================================================
\section{Coffee shop: the initial cash flow --- Berk 9.18, p.~276}

\begin{notation}
  \item $S=20{,}000$: feasibility study already paid (sunk). $P=100{,}000$: price paid for the space two years ago (sunk).
  \item $L=115{,}000$: after-tax amount the space would fetch if sold today (opportunity cost).
  \item $K=30{,}000$: capital expenditure on outfitting; $N=5{,}000$: inventory (investment in net working capital).
  \item $\mathrm{CF}_0$: incremental cash flow at $t=0$, in dollars.
\end{notation}

\step{1}{Incremental principle} Only cash flows that change because the shop is opened count.
$S$ and $P$ have already been paid whatever is decided, so they are excluded. Using the space
for the shop means giving up the sale of it, so the after-tax sale value $L$ is a cost of the
project:
\begin{equation}\label{eq:3cf}
  \mathrm{CF}_0 = -(L + K + N).
\end{equation}

\step{2}{Substitute} Placing the numbers in \eqref{eq:3cf}:
\begin{align}
  \mathrm{CF}_0 &= -\big((115{,}000) + (30{,}000) + (5{,}000)\big) \notag\\
                &= -150{,}000. \label{eq:3res}
\end{align}

\intuition{The question is ``what do I give up today to run the shop?''. The answer is the
\$115{,}000 I could pocket by selling the space, plus the new money for equipment and stock.
The study and the old purchase price are gone either way.}

\refline{Berk, \S9.3--9.4 (incremental free cash flow; opportunity costs, sunk costs), pp.~254--265.}

% =============================================================================
\section{Pre-tax WACC and the return on assets --- Berk 13.4, p.~403}

\begin{notation}
  \item $A_0=1000$, $D_0=400$, $E_0=600$: market values today of assets, debt and equity, $A_0=D_0+E_0$.
  \item $A_1^H=1200$, $A_1^L=960$: asset values next year, each with probability $1/2$.
  \item $r_D=0.05$: promised interest rate on debt; $D_1$: payment to debt next year.
  \item $E_1^H$, $E_1^L$: equity payoffs next year; $r_E$, $r_A$: expected returns on equity and assets.
  \item $w_D=D_0/A_0=0.4$, $w_E=E_0/A_0=0.6$: capital-structure weights.
\end{notation}

\step{1}{Debt payoff} Debt is promised $D_0(1+r_D)$:
\begin{align}
  D_1 &= (400)(1.05) \notag\\
      &= 420. \label{eq:4d1}
\end{align}
Since $A_1^L=960>420$, the promise is paid in both states, so debt is riskless and its
expected return is $r_D=0.05$.

\step{2}{Equity payoffs} Equity receives what is left, $E_1^s=A_1^s-D_1$. Substituting \eqref{eq:4d1}:
\begin{align}
  E_1^H &= 1200-(420) = 780, \notag\\
  E_1^L &= 960-(420) = 540. \label{eq:4e1}
\end{align}

\step{3}{Expected equity return} Using \eqref{eq:4e1}:
\begin{align}
  E[E_1] &= (0.5)(780) + (0.5)(540) \notag\\
         &= 390 + 270 \notag\\
         &= 660, \notag\\
  r_E &= \frac{(660)}{(600)} - 1 \notag\\
      &= 1.10 - 1 \notag\\
      &= 0.10. \label{eq:4re}
\end{align}

\step{4}{Expected asset return}
\begin{align}
  E[A_1] &= (0.5)(1200) + (0.5)(960) \notag\\
         &= 600 + 480 \notag\\
         &= 1080, \notag\\
  r_A &= \frac{(1080)}{(1000)} - 1 \notag\\
      &= 0.08. \label{eq:4ra}
\end{align}

\step{5}{Return on the 40/60 portfolio (pre-tax WACC)} Substituting $r_D$ and \eqref{eq:4re}:
\begin{align}
  w_D r_D + w_E r_E &= (0.4)(0.05) + (0.6)(0.10) \notag\\
                    &= 0.02 + 0.06 \notag\\
                    &= 0.08, \label{eq:4wacc}
\end{align}
which equals $r_A$ in \eqref{eq:4ra}.

\step{6}{Why it must hold (general proof)} With $w_D=D_0/A_0$ and $w_E=E_0/A_0$,
\begin{align}
  w_D r_D + w_E r_E &= \frac{D_0}{A_0}\Big(\frac{E[D_1]}{D_0}-1\Big) + \frac{E_0}{A_0}\Big(\frac{E[E_1]}{E_0}-1\Big) \notag\\
  &= \frac{E[D_1]+E[E_1]}{A_0} - \frac{D_0+E_0}{A_0} \notag\\
  &= \frac{E[A_1]}{A_0} - 1 = r_A, \label{eq:4gen}
\end{align}
where the last line uses $D_1+E_1=A_1$ state by state (all asset cash flows go to debt or
equity) and $D_0+E_0=A_0$.

\intuition{Holding 40\% of the debt and 60\% of the equity is holding 100\% of the claims on the
assets, scaled down. A portfolio that owns the whole cash flow of the assets must earn what the
assets earn; leverage only splits that return between the two claims.}

\refline{Berk, \S13.1 (a first look at the WACC), pp.~383--386; \S13.3, p.~390.}

% =============================================================================
\section{WACC for SKYE Corporation --- Higgins 8.7, p.~328}

\begin{notation}
  \item $E$: market value of equity (millions of dollars); $D$: value of interest-bearing debt (millions), $D,E>0$.
  \item $K_D=0.0768$: pre-tax cost of debt (yield to maturity); $t=0.40$: tax rate, $0\le t<1$.
  \item $r_f=0.06$: government bond rate; $\beta=0.75$: equity beta; $\mathrm{MRP}=0.063$: historical excess return on stocks.
  \item $K_E$: cost of equity; $K_W$: weighted-average cost of capital.
\end{notation}

\step{1}{Weights} Equity at market value, $E=(30)(1)=30$. The market value of the debt is not
given (no maturity), so, as Higgins does, use its book value $D=10$. Total $D+E=40$.

\step{2}{Cost of equity (CAPM)}
\begin{align}
  K_E &= r_f + \beta\,\mathrm{MRP} \label{eq:5ke}\\
      &= (0.06) + (0.75)(0.063) \notag\\
      &= 0.06 + 0.04725 \notag\\
      &= 0.10725. \label{eq:5keres}
\end{align}

\step{3}{After-tax cost of debt} The relevant rate is the yield, not the 9\% coupon:
\begin{align}
  (1-t)K_D &= (1-0.40)(0.0768) \notag\\
           &= 0.04608. \label{eq:5kd}
\end{align}

\step{4}{WACC} Higgins' Equation 8.1, $K_W = \dfrac{(1-t)K_D D + K_E E}{D+E}$. Substituting
\eqref{eq:5kd}, \eqref{eq:5keres} and Step 1:
\begin{align}
  K_W &= \frac{(0.04608)(10) + (0.10725)(30)}{(40)} \notag\\
      &= \frac{0.4608 + 3.2175}{40} \notag\\
      &= \frac{3.6783}{40} \notag\\
      &= 0.0919575 \approx 9.2\%. \label{eq:5res}
\end{align}

\intuition{The firm's capital is three-quarters equity, so the WACC sits close to the 10.7\% cost
of equity, pulled down by cheap debt whose cost is cut further by the interest tax deduction.}

\remark{The 9\% coupon above the 7.68\% yield means the bonds trade above par, so the market
value of debt exceeds \$10 million and the true market-weight WACC is slightly lower. Without a
maturity the market value cannot be computed; book value is the stated fallback.}

\bookcheck{Matches the book's answer, $K_W=9.2\%$ (p.~422).}

\refline{Higgins, \S8.4.1--8.4.2 (cost of capital defined; cost of capital for Polaris, CAPM), pp.~295--304.}

% =============================================================================
\section{GM--Hughes: cost of capital from comparables --- Grinblatt 11.10--11.13, p.~416}

\begin{notation}
  \item $\beta_E$: equity beta; $D/E$: debt-to-equity ratio (market values); $\beta_A$: asset (unlevered) beta.
  \item $\beta_D=0$: debt beta (debt is riskless, discounted at $r_f$). No taxes.
  \item $r_f=0.08$: riskless rate; $\bar R_T=0.14$: expected return of the tangency portfolio.
  \item $r_A$: cost of capital of Hughes' assets; $C_1=300$: expected cash flow next year (\$ millions); $g=0.05$: growth rate, $g<r_A$.
  \item $V_H$: value of Hughes (\$ millions).
\end{notation}

\step{1}{Asset beta of a levered firm (11.10)} Assets are a portfolio of the firm's debt and
equity, so their beta is the value-weighted average:
\begin{equation}\label{eq:6ba}
  \beta_A = \frac{E}{D+E}\,\beta_E + \frac{D}{D+E}\,\beta_D.
\end{equation}
Substituting $\beta_D=0$ and dividing numerator and denominator by $E$:
\begin{equation}\label{eq:6ba2}
  \beta_A = \frac{\beta_E}{1+D/E}.
\end{equation}
Lockheed, substituting into \eqref{eq:6ba2}:
\begin{align}
  \beta_{A,L} &= \frac{(0.90)}{1+(0.90)} = \frac{0.90}{1.90} = 0.4736842. \label{eq:6bl}
\end{align}
Northrop, substituting into \eqref{eq:6ba2}:
\begin{align}
  \beta_{A,N} &= \frac{(0.85)}{1+(0.70)} = \frac{0.85}{1.70} = 0.5000000. \label{eq:6bn}
\end{align}

\step{2}{Hughes asset beta (11.11)} Averaging \eqref{eq:6bl} and \eqref{eq:6bn}:
\begin{align}
  \beta_{A,H} &= \frac{(0.4736842) + (0.5000000)}{2} \notag\\
              &= \frac{0.9736842}{2} \notag\\
              &= 0.4868421. \label{eq:6bh}
\end{align}

\step{3}{Cost of capital (11.12)} With the CAPM ($\bar R_T$ plays the role of the market return),
substituting \eqref{eq:6bh}:
\begin{align}
  r_A &= r_f + \beta_{A,H}(\bar R_T - r_f) \label{eq:6ra}\\
      &= (0.08) + (0.4868421)(0.14-0.08) \notag\\
      &= 0.08 + 0.0292105 \notag\\
      &= 0.1092105 \approx 10.92\%. \label{eq:6rares}
\end{align}
The target $D/E=1$ does not change this: without taxes the cost of the assets is independent of
how they are financed (MM). Check: relevering \eqref{eq:6ba2} at $D/E=1$ gives
$\beta_{E,H}=(0.4868421)(1+1)=0.9736842$, so $r_E=(0.08)+(0.9736842)(0.06)=0.08+0.0584211=0.1384211$
and
\begin{align}
  \mathrm{WACC} &= (0.5)(0.08) + (0.5)(0.1384211) \notag\\
                &= 0.04 + 0.0692105 \notag\\
                &= 0.1092105, \notag
\end{align}
the same as \eqref{eq:6rares}.

\step{4}{Value of Hughes (11.13)} The cash flow grows at $g$ forever, so it is a growing perpetuity.
Substituting \eqref{eq:6rares}:
\begin{align}
  V_H &= \frac{C_1}{r_A-g} \label{eq:6v}\\
      &= \frac{(300)}{(0.1092105)-(0.05)} \notag\\
      &= \frac{300}{0.0592105} \notag\\
      &= 5{,}066.67 \text{ (\$ millions)}. \label{eq:6vres}
\end{align}

\remark{Using GM's own beta would be wrong: $\beta_{A,GM}=1.20/1.40=0.8571429$, giving
$r=(0.08)+(0.8571429)(0.06)=0.08+0.0514286=0.1314286$ and $V=300/(0.1314286-0.05)=300/0.0814286=3{,}684.21$,
undervaluing Hughes by about \$1.4 billion.}

\intuition{The discount rate belongs to the project, not to the buyer. Hughes' risk is read off
firms in the same business, after stripping out their leverage, because leverage scales equity
risk without changing the risk of the operations.}

\refline{Grinblatt \& Titman, \S11.3 (the effect of leverage on comparisons), p.~379; \S11.4 and
Example~11.2 (comparison method, adjusting for leverage without taxes), pp.~384--386.}

% =============================================================================
\section{Mutually exclusive projects and the crossover rate --- Berk 8.26, p.~244}

\begin{notation}
  \item $C^A_t$, $C^B_t$: cash flows of projects $A$ and $B$ at $t=0,\dots,4$ (as tabulated in the statement).
  \item $r$: discount rate, $r>-1$; $r_0=0.11$: the cost of capital.
  \item $\mathrm{NPV}_j(r)=\sum_{t=0}^{4} C^j_t/(1+r)^t$, $j\in\{A,B\}$; $\mathrm{IRR}_j$: root of $\mathrm{NPV}_j(r)=0$.
  \item $\Delta_t=C^A_t-C^B_t$: incremental cash flows of $A$ over $B$; $r^\ast$: crossover rate.
\end{notation}

\step{1}{(a) NPVs at 11\%} Substituting $r_0$ in the definition:
\begin{align}
  \mathrm{NPV}_A &= -100 + \frac{25}{(1.11)} + \frac{30}{(1.11)^2} + \frac{40}{(1.11)^3} + \frac{50}{(1.11)^4} \notag\\
                 &= -100 + 22.5225 + 24.3487 + 29.2477 + 32.9365 \notag\\
                 &= 9.055, \label{eq:7na}\\
  \mathrm{NPV}_B &= -100 + \frac{50}{(1.11)} + \frac{40}{(1.11)^2} + \frac{30}{(1.11)^3} + \frac{20}{(1.11)^4} \notag\\
                 &= -100 + 45.0450 + 32.4649 + 21.9357 + 13.1746 \notag\\
                 &= 12.620. \label{eq:7nb}
\end{align}

\step{2}{(b) IRRs} Both streams have one sign change ($-$ then $+$), so $\mathrm{NPV}_j(r)$ is
strictly decreasing in $r$ and has exactly one root. Solving $\mathrm{NPV}_j(r)=0$ numerically:
\begin{equation}\label{eq:7irr}
  \mathrm{IRR}_A = 14.70\%, \qquad \mathrm{IRR}_B = 17.80\%.
\end{equation}

\step{3}{(c) Crossover rate} Indifference means $\mathrm{NPV}_A(r)=\mathrm{NPV}_B(r)$, i.e.\
the incremental project $\Delta$ has zero NPV:
\begin{align}
  \Delta &= (0-0,\ 25-50,\ 30-40,\ 40-30,\ 50-20) = (0,\,-25,\,-10,\,10,\,30), \notag\\
  0 &= \frac{-25}{1+r} + \frac{-10}{(1+r)^2} + \frac{10}{(1+r)^3} + \frac{30}{(1+r)^4}. \label{eq:7cross}
\end{align}
Multiplying \eqref{eq:7cross} by $(1+r)^4$ and writing $x=1+r>0$:
\begin{equation}\label{eq:7poly}
  -25x^3 - 10x^2 + 10x + 30 = 0.
\end{equation}
The coefficients change sign once, so by Descartes' rule \eqref{eq:7poly} has exactly one
positive root. Solving numerically,
\begin{equation}\label{eq:7rstar}
  r^\ast = 5.57\%, \qquad \mathrm{NPV}_A(r^\ast) = \mathrm{NPV}_B(r^\ast) = 24.86.
\end{equation}
For $r<r^\ast$, $A$ has the higher NPV (its cash arrives later, so it gains more from a low
rate); for $r>r^\ast$, $B$ does.

\step{4}{(d) Decision} At $r_0=11\%>r^\ast$, \eqref{eq:7nb} exceeds \eqref{eq:7na}: take project
$B$. Both IRRs in \eqref{eq:7irr} exceed 11\%, so each project alone is acceptable; the choice
between them follows the NPV. The IRR ranking happens to agree here, but it would rank wrongly
at any cost of capital below $5.57\%$.

\intuition{IRR measures a rate, NPV measures value. When projects differ in timing, the one with
earlier cash flows can have the higher IRR yet the lower NPV at low discount rates. The
crossover rate is exactly the IRR of the difference between the projects.}

\refline{Berk, \S8.3 (alternative rules, IRR), pp.~216--226; \S8.4 (choosing between projects,
crossover), pp.~226--232.}

% =============================================================================
\section{Innovation Company: how many IRRs? --- Berk 8.17, p.~243}

\begin{notation}
  \item $I_0=5$: up-front cost (\$ millions); $\pi=1$: profit per year for years 1--10 (\$ millions).
  \item $s=0.1$: support cost per year in perpetuity, from year 1 (\$ millions).
  \item $r>0$: cost of capital; $a(r)=\big(1-(1+r)^{-10}\big)/r$: ten-year annuity factor.
  \item $N(r)$: NPV of the project at rate $r$ (\$ millions).
\end{notation}

\step{1}{NPV function} An annuity of profits plus a perpetuity of costs:
\begin{equation}\label{eq:8n}
  N(r) = -I_0 + \pi\,a(r) - \frac{s}{r}.
\end{equation}

\step{2}{(a) Evaluate at 6\%, 2\% and 11\%} Substituting each rate into \eqref{eq:8n}:
\begin{align}
  N(0.06) &= -5 + (1)(7.36009) - \frac{0.1}{0.06} \notag\\
          &= -5 + 7.36009 - 1.66667 \notag\\
          &= 0.69342, \label{eq:8n6}\\
  N(0.02) &= -5 + (1)(8.98259) - \frac{0.1}{0.02} \notag\\
          &= -5 + 8.98259 - 5.00000 \notag\\
          &= -1.01741, \label{eq:8n2}\\
  N(0.11) &= -5 + (1)(5.88923) - \frac{0.1}{0.11} \notag\\
          &= -5 + 5.88923 - 0.90909 \notag\\
          &= -0.01986. \label{eq:8n11}
\end{align}
Undertake the project at 6\% (NPV $\approx$ \$693{,}420); reject it at 2\% and at 11\%.

\step{3}{(b) Number of IRRs} For $r>0$, multiply \eqref{eq:8n} by $r$ (which preserves the
zeros):
\begin{align}
  r\,N(r) &= -5r + \big(1-(1+r)^{-10}\big) - 0.1 \notag\\
          &= 0.9 - (1+r)^{-10} - 5r. \notag
\end{align}
So $N(r)=0$ if and only if
\begin{equation}\label{eq:8g}
  g(r)\equiv 0.9-(1+r)^{-10} = 5r.
\end{equation}
Since $(1+r)^{-10}$ is strictly convex, $g$ is strictly concave, and a strictly concave function
crosses a straight line at most twice. By \eqref{eq:8n2}--\eqref{eq:8n11}, $N$ changes sign on
$(0.02,0.06)$ and again on $(0.06,0.11)$, so there are exactly two IRRs. Solving numerically:
\begin{equation}\label{eq:8irr}
  \mathrm{IRR}_1 = 2.75\%, \qquad \mathrm{IRR}_2 = 10.88\%.
\end{equation}

\step{4}{(c) What the IRR rule says} The rule ``accept if the IRR exceeds the cost of capital'' is
not defined when there are two IRRs: at $r=6\%$ one IRR is above $r$ and the other below.
The NPV is positive exactly for $2.75\%<r<10.88\%$ (by \eqref{eq:8g}, between the two
crossings) and negative outside, so only the NPV rule gives the right answer. At very low rates
the perpetual support cost, $s/r$, dominates; at high rates the up-front cost does.

\intuition{The cash flows change sign twice ($-$, then $+$ for ten years, then $-$ forever), so
the project is partly an investment and partly a loan that is never repaid. Different rates
favour different parts, and NPV is positive only in a window of rates.}

\refline{Berk, \S8.3 (IRR rule pitfalls: multiple IRRs), pp.~216--226.}

% =============================================================================
\section{CoffeeStop: WACC for the liquor division --- Berk 13.20, p.~405}

\begin{notation}
  \item $\beta_E=0.26$: equity beta of the pure-play comparable BF Liquors; $E/V=0.89$, $D/V=0.11$: its weights.
  \item $r_f=0.03$; $\mathrm{MRP}=0.06$: market risk premium; $r_D=0.048$: cost of debt; $\tau=0.35$: tax rate.
  \item $r_E$: cost of equity; $r_U$: unlevered (pre-tax) cost of capital; $r_{\mathrm{wacc}}$: after-tax WACC.
\end{notation}

\step{1}{Choose the comparable} The division sells liquor, so CoffeeStop's own beta (coffee) is
the wrong risk. BF Liquors is the pure play, and it has the same leverage, $D/V=0.11$, that is
planned for the division. Its WACC is therefore the division's WACC (Berk, divisional costs of
capital).

\step{2}{Cost of equity (CAPM)}
\begin{align}
  r_E &= (0.03) + (0.26)(0.06) \notag\\
      &= 0.03 + 0.0156 \notag\\
      &= 0.0456. \label{eq:9re}
\end{align}

\step{3}{WACC} Substituting \eqref{eq:9re} into $r_{\mathrm{wacc}} = \frac{E}{V}r_E + \frac{D}{V}r_D(1-\tau)$:
\begin{align}
  r_{\mathrm{wacc}} &= (0.89)(0.0456) + (0.11)(0.048)(1-0.35) \notag\\
                    &= 0.040584 + 0.003432 \notag\\
                    &= 0.044016 \approx 4.40\%. \label{eq:9wacc}
\end{align}

\step{4}{Check: unlever and relever} The pre-tax WACC equals the unlevered cost of capital (as
proved in Exercise 4):
\begin{align}
  r_U &= (0.89)(0.0456) + (0.11)(0.048) \notag\\
      &= 0.040584 + 0.00528 \notag\\
      &= 0.045864. \label{eq:9ru}
\end{align}
The after-tax WACC subtracts the interest tax shield, $r_{\mathrm{wacc}} = r_U - \frac{D}{V}\tau r_D$,
which follows from the two definitions. Substituting \eqref{eq:9ru} at the division's $D/V=0.11$:
\begin{align}
  r_{\mathrm{wacc}} &= (0.045864) - (0.11)(0.35)(0.048) \notag\\
                    &= 0.045864 - 0.001848 \notag\\
                    &= 0.044016, \notag
\end{align}
the same as \eqref{eq:9wacc}, as it must be when the division copies the comparable's leverage.

\remark{The data are not fully consistent under the CAPM: the cost of debt implies a debt beta
$(0.048-0.03)/0.06=0.30$, above the equity beta of $0.26$, so $r_E<r_D$. The answer uses the
inputs as given. Using CoffeeStop's own WACC, $(0.96)(0.03+0.61\cdot0.06)+(0.04)(0.048)(0.65)
=(0.96)(0.0666)+0.001248=0.065184$, would overstate the division's rate by about 2.1 points.}

\intuition{A firm should discount each line of business at the rate investors demand for that
business. A liquor division is valued like a liquor company, not like the parent.}

\refline{Berk, \S13.3 (WACC), p.~390; \S13.5 (project-based costs of capital, divisional costs),
pp.~396--398.}

% =============================================================================
\section{Karamazov Semiconductors: enterprise vs.\ equity IRR --- Higgins 8.16, p.~330}

\begin{notation}
  \item $I=10$: investment (\$ millions); $C=1$: after-tax (unlevered) cash flow per year in perpetuity (\$ millions).
  \item $D/V=E/V=0.5$: financing weights; $D=5$, $E=5$: debt and equity raised (\$ millions).
  \item $K_D=0.08$: pre-tax cost of debt; $t=0.25$: tax rate; $K_E=0.14$: cost of equity; $K_W$: WACC.
  \item Assumption: the debt is perpetual (kept at 50\% of the project), interest paid annually.
\end{notation}

\step{1}{(a) WACC}
\begin{align}
  K_W &= \frac{D}{V}(1-t)K_D + \frac{E}{V}K_E \notag\\
      &= (0.5)(1-0.25)(0.08) + (0.5)(0.14) \notag\\
      &= 0.03 + 0.07 \notag\\
      &= 0.10. \label{eq:10kw}
\end{align}

\step{2}{(b) Enterprise IRR} Enterprise cash flows are $-I$ then $C$ forever. The IRR solves
$-I + C/\mathrm{IRR}=0$:
\begin{align}
  \mathrm{IRR}_{\mathrm{ent}} &= \frac{C}{I} = \frac{(1)}{(10)} \notag\\
                              &= 0.10. \label{eq:10ent}
\end{align}

\step{3}{(c) Equity IRR} Equity puts up $E=5$ and receives $C$ minus after-tax interest:
\begin{align}
  K_D(1-t)D &= (0.08)(1-0.25)(5) \notag\\
            &= 0.3, \label{eq:10int}\\
  C_E &= C - K_D(1-t)D = (1) - (0.3) \notag\\
      &= 0.7, \label{eq:10ce}\\
  \mathrm{IRR}_{\mathrm{eq}} &= \frac{C_E}{E} = \frac{(0.7)}{(5)} \notag\\
                             &= 0.14. \label{eq:10eq}
\end{align}

\step{4}{(d) Why they differ} The two IRRs measure returns on different cash flows: the
enterprise IRR \eqref{eq:10ent} is the return on all capital, the equity IRR \eqref{eq:10eq} is
the return on the shareholders' part only, after paying debt holders. Leverage concentrates the
surplus over the after-tax cost of debt on a smaller equity base, so the equity IRR is higher
and also riskier. Each IRR must be compared with its own hurdle: $10\%$ against $K_W=10\%$ and
$14\%$ against $K_E=14\%$. Both comparisons give the same verdict: the project earns exactly
its cost of capital (NPV $=0$). Comparing the equity IRR with the WACC (or the enterprise IRR
with the cost of equity) would mix perspectives.

\intuition{Debt does not create the extra 4 points of equity return; it transfers risk. Viewed
consistently, the enterprise and equity perspectives always agree on whether a project creates
value.}

\refline{Higgins, \S8.6.1 (the enterprise perspective vs.\ the equity perspective), pp.~307--310.}

% =============================================================================
\section{Percolated Fiber: fixing the consultants' report --- Berk 9.25, p.~277}

\begin{notation}
  \item Amounts in thousands of dollars; $t=0,\dots,10$ years.
  \item $K=25{,}000$: capital expenditure at $t=0$; $\mathrm{Dep}=2{,}500$: straight-line depreciation, years 1--10.
  \item $\mathrm{NWC}=10{,}000$: working capital invested at $t=0$, recovered at $t=10$.
  \item $\mathrm{SGA}=1{,}000$: incremental SG\&A (the other 1{,}000 is overhead incurred anyway).
  \item $\tau$: tax rate; $\mathrm{FCF}_t$: free cash flow; $r=0.14$: cost of capital; $a$: ten-year annuity factor at 14\%.
\end{notation}

\step{1}{Tax rate} From the report, $\tau = 2{,}625/7{,}500 = 0.35$.

\step{2}{Irrelevant items} The \$1 million consultants' fee is owed whatever is decided (sunk):
excluded. Only the incremental 1{,}000 of SG\&A enters.

\step{3}{Incremental EBIT and unlevered net income (years 1--10)}
\begin{align}
  \mathrm{EBIT} &= 30{,}000 - 18{,}000 - (1{,}000) - 2{,}500 \notag\\
                &= 8{,}500, \notag\\
  \text{Tax} &= (0.35)(8{,}500) = 2{,}975, \notag\\
  \text{Unlevered NI} &= 8{,}500 - (2{,}975) = 5{,}525. \label{eq:11ni}
\end{align}

\step{4}{(a) Free cash flows} $\mathrm{FCF}_t = \text{Unlevered NI}_t + \mathrm{Dep}_t - \mathrm{CapEx}_t - \Delta\mathrm{NWC}_t$:
\begin{align}
  \mathrm{FCF}_0 &= -(25{,}000) - (10{,}000) = -35{,}000, \notag\\
  \mathrm{FCF}_{1\ldots9} &= (5{,}525) + (2{,}500) = 8{,}025, \label{eq:11fcf}\\
  \mathrm{FCF}_{10} &= (8{,}025) + (10{,}000) = 18{,}025. \notag
\end{align}

\step{5}{(b) Value at 14\%} Writing $\mathrm{FCF}_{10}$ as $8{,}025$ plus the NWC recovery:
\begin{align}
  \mathrm{NPV} &= -35{,}000 + 8{,}025\,a + \frac{10{,}000}{(1.14)^{10}} \notag\\
               &= -35{,}000 + (8{,}025)(5.216116) + \frac{10{,}000}{(3.707221)} \notag\\
               &= -35{,}000 + 41{,}859.33 + 2{,}697.44 \notag\\
               &= 9{,}556.77. \label{eq:11npv}
\end{align}
The project adds about \$9.56 million of value (the future cash flows are worth \$44.56 million
against \$35 million invested), far from the consultants' \$48.75 million.

\intuition{Earnings are not cash: the equipment is paid for up front, not in yearly instalments
of depreciation; working capital ties up cash until year 10; and money in year 10 is worth much
less than money today. Overhead that exists anyway is not caused by the project.}

\refline{Berk, \S9.2--9.3 (incremental earnings and free cash flow), pp.~250--260; \S9.4 (NWC,
sunk costs, overhead), pp.~260--265.}

% =============================================================================
\section{Torino Marine: three terminal values --- Higgins 9.11, p.~382}

\begin{notation}
  \item Amounts in \$ millions; valuation date end of 2021; $\mathrm{FCF}_{2022},\dots,\mathrm{FCF}_{2025}=-52,76,92,112$.
  \item $K_W=0.11$: WACC; $t=0.40$: tax rate; $\mathrm{EBIT}$: earnings before interest and tax.
  \item $TV_{2025}$: terminal (enterprise) value at end of 2025; $\mathrm{FMV}$: enterprise value at end of 2021.
  \item $D=250$: interest-bearing liabilities; $n=40$: shares (millions); $g=0.05$: growth in (c).
\end{notation}

\step{1}{PV of the explicit forecast (common to all parts)}
\begin{align}
  PV_{FCF} &= \frac{-52}{(1.11)} + \frac{76}{(1.11)^2} + \frac{92}{(1.11)^3} + \frac{112}{(1.11)^4} \notag\\
           &= -46.847 + 61.683 + 67.270 + 73.778 \notag\\
           &= 155.884. \label{eq:12pv}
\end{align}
The discount factor to 2025 is $(1.11)^4 = 1.518070$.

\step{2}{(a) No-growth terminal value} With depreciation equal to capex and no change in
working capital, $\mathrm{FCF}=\mathrm{EBIT}(1-t)$ from 2026 on, a level perpetuity:
\begin{align}
  TV_{2025} &= \frac{\mathrm{EBIT}(1-t)}{K_W} = \frac{(200)(1-0.40)}{(0.11)} \notag\\
            &= \frac{120}{0.11} = 1{,}090.909, \notag\\
  PV(TV) &= \frac{1{,}090.909}{(1.518070)} = 718.616, \notag\\
  \mathrm{FMV}_a &= (155.884) + (718.616) = 874.500. \label{eq:12a}
\end{align}

\step{3}{(b) Equity per share} Equity is enterprise value minus debt. Substituting \eqref{eq:12a}:
\begin{align}
  \frac{\mathrm{FMV}_a - D}{n} &= \frac{(874.500) - (250)}{(40)} \notag\\
                               &= \frac{624.500}{40} = 15.61 \text{ dollars per share}. \label{eq:12b}
\end{align}

\step{4}{(c) Growing terminal value} 2026 EBIT is $200(1.05)$; FCF subtracts net capex (30) and
working-capital investment (15):
\begin{align}
  \mathrm{FCF}_{2026} &= (200)(1.05)(1-0.40) - 30 - 15 \notag\\
                      &= (210)(0.60) - 45 \notag\\
                      &= 126 - 45 = 81, \notag\\
  TV_{2025} &= \frac{(81)}{(0.11)-(0.05)} = \frac{81}{0.06} = 1{,}350, \notag\\
  PV(TV) &= \frac{1{,}350}{(1.518070)} = 889.287, \notag\\
  \mathrm{FMV}_c &= (155.884) + (889.287) = 1{,}045.171, \notag\\
  \text{per share} &= \frac{(1{,}045.171) - (250)}{(40)} = \frac{795.171}{40} = 19.88. \label{eq:12c}
\end{align}

\step{5}{(d) Multiple-based terminal value} 2025 net income deducts interest on the debt:
\begin{align}
  NI_{2025} &= \big(200 - (0.10)(250)\big)(1-0.40) \notag\\
            &= (175)(0.60) = 105, \notag\\
  E_{2025} &= (12)(105) = 1{,}260, \notag\\
  TV_{2025} &= E_{2025} + D = (1{,}260) + (250) = 1{,}510, \notag\\
  PV(TV) &= \frac{1{,}510}{(1.518070)} = 994.684, \notag\\
  \mathrm{FMV}_d &= (155.884) + (994.684) = 1{,}150.568, \notag\\
  \text{per share} &= \frac{(1{,}150.568) - (250)}{(40)} = \frac{900.568}{40} = 22.51. \label{eq:12d}
\end{align}
The P/E multiple values equity, so debt is added back to get an enterprise terminal value that
can be discounted at the WACC together with the FCFs.

\intuition{Most of the value (82--86\%) sits in the terminal value, so the method chosen for it
drives the answer: \$15.61, \$19.88 or \$22.51 per share. Growth is worth something only because
the reinvestment it needs (\$45 million in 2026) earns more than the WACC; the market multiple
of 12 implies even more optimism than 5\% growth.}

\bookcheck{Matches the book's answers (p.~427): (a) \$874.5 million, (b) \$15.61, (c)
\$1{,}045.2 million and \$19.88, (d) \$22.51 per share. In (d) the book adds \$155.8 instead of
\$155.9 for the forecast-period PV and prints \$1{,}150.5 million; carried to three decimals the
enterprise value is \$1{,}150.568 million. The per-share value is unaffected.}

\refline{Higgins, \S9.2.1--9.2.3 (free cash flow, the terminal value, a numerical example),
pp.~342--348; \S9.3 (valuation by comparables, P/E), p.~349.}

% =============================================================================
\section{Calculating project NPV (CUTO) --- Ross 10.31 (Desafio), p.~356}

\begin{notation}
  \item Amounts in dollars; $t=0,\dots,4$ years; $\tau=0.34$: tax rate; $R=0.13$: required return.
  \item $L_0=1{,}500{,}000$: after-tax value of the land today (opportunity cost); $L_4=1{,}600{,}000$: after-tax land value at $t=4$.
  \item $q_t$: units sold (3{,}200, 4{,}300, 3{,}900, 2{,}800); $P=780$: price; $S_t=Pq_t$: sales.
  \item $\mathrm{VC}_t=0.15\,S_t$: variable cost; $\mathrm{FC}=425{,}000$: fixed cost per year.
  \item $K=4{,}200{,}000$: equipment; $\mathrm{Dep}_t$: tax depreciation; $SV=400{,}000$: salvage price at $t=4$.
  \item $\mathrm{NWC}=125{,}000$: working capital at $t=0$; $\mathrm{OCF}_t$: operating cash flow; $\mathrm{CF}_t$: total cash flow.
\end{notation}

\step{1}{Relevant and irrelevant items} The \$1.4 million paid for the land three years ago and
the \$125{,}000 marketing study are sunk: excluded. Using the land forgoes selling it now for
$L_0$, so $L_0$ is a $t=0$ outflow; at the end the land can still be sold for $L_4$, a $t=4$
inflow. ``Assume the firm has other profitable projects'' means any tax loss could be used
against other income; it does not bind here because EBIT is positive in every year
(Step~3).

\step{2}{Depreciation} The equipment is ``\pt{depreciado de forma linear e acelerada, em tr\^es
anos}'': straight line to zero over three years, as in Ross's Problem~9 of the same chapter (p.~353):
\begin{align}
  \mathrm{Dep}_{1,2,3} &= \frac{(4{,}200{,}000)}{3} = 1{,}400{,}000, \qquad \mathrm{Dep}_4 = 0. \label{eq:13dep}
\end{align}
The book value at $t=4$ is therefore zero.

\step{3}{Operating cash flow} For each year, $\mathrm{EBIT}_t = S_t - \mathrm{VC}_t - \mathrm{FC} - \mathrm{Dep}_t$
and $\mathrm{OCF}_t = \mathrm{EBIT}_t(1-\tau) + \mathrm{Dep}_t$. Year 1 in full:
\begin{align}
  S_1 &= (3{,}200)(780) = 2{,}496{,}000, \notag\\
  \mathrm{VC}_1 &= (0.15)(2{,}496{,}000) = 374{,}400, \notag\\
  \mathrm{EBIT}_1 &= 2{,}496{,}000 - (374{,}400) - (425{,}000) - (1{,}400{,}000) = 296{,}600, \notag\\
  \text{Tax}_1 &= (0.34)(296{,}600) = 100{,}844, \notag\\
  NI_1 &= 296{,}600 - (100{,}844) = 195{,}756, \notag\\
  \mathrm{OCF}_1 &= (195{,}756) + (1{,}400{,}000) = 1{,}595{,}756. \label{eq:13ocf1}
\end{align}
The other years follow the same lines:
\begin{center}
\begin{tabular}{lrrrr}
\toprule
Year & 1 & 2 & 3 & 4\\
\midrule
Sales $S_t$          & 2{,}496{,}000 & 3{,}354{,}000 & 3{,}042{,}000 & 2{,}184{,}000\\
Variable cost (15\%) & 374{,}400 & 503{,}100 & 456{,}300 & 327{,}600\\
Fixed cost           & 425{,}000 & 425{,}000 & 425{,}000 & 425{,}000\\
Depreciation         & 1{,}400{,}000 & 1{,}400{,}000 & 1{,}400{,}000 & 0\\
EBIT                 & 296{,}600 & 1{,}025{,}900 & 760{,}700 & 1{,}431{,}400\\
Tax (34\%)           & 100{,}844 & 348{,}806 & 258{,}638 & 486{,}676\\
Net income           & 195{,}756 & 677{,}094 & 502{,}062 & 944{,}724\\
OCF                  & 1{,}595{,}756 & 2{,}077{,}094 & 1{,}902{,}062 & 944{,}724\\
\bottomrule
\end{tabular}
\end{center}

\step{4}{After-tax salvage} The equipment is sold for $SV$ with book value zero, so the whole
price is a taxable gain:
\begin{align}
  SV(1-\tau) &= 400{,}000 - (0.34)(400{,}000 - 0) \notag\\
             &= 400{,}000 - 136{,}000 = 264{,}000. \label{eq:13sv}
\end{align}

\step{5}{Total cash flows} Assuming the working capital is fully recovered when sales stop:
\begin{align}
  \mathrm{CF}_0 &= -(4{,}200{,}000) - (1{,}500{,}000) - (125{,}000) = -5{,}825{,}000, \notag\\
  \mathrm{CF}_{1,2,3} &= \mathrm{OCF}_{1,2,3} = 1{,}595{,}756,\ 2{,}077{,}094,\ 1{,}902{,}062, \notag\\
  \mathrm{CF}_4 &= (944{,}724) + (1{,}600{,}000) + (264{,}000) + (125{,}000) = 2{,}933{,}724. \label{eq:13cf}
\end{align}

\step{6}{NPV at 13\%} Discounting \eqref{eq:13cf}:
\begin{align}
  \mathrm{NPV} &= -5{,}825{,}000 + \frac{1{,}595{,}756}{(1.13)} + \frac{2{,}077{,}094}{(1.13)^2}
                  + \frac{1{,}902{,}062}{(1.13)^3} + \frac{2{,}933{,}724}{(1.13)^4} \notag\\
               &= -5{,}825{,}000 + 1{,}412{,}173.45 + 1{,}626{,}669.28 + 1{,}318{,}224.38 + 1{,}799{,}307.87 \notag\\
               &= 331{,}374.98. \label{eq:13npv}
\end{align}
The NPV is positive: CUTO should go ahead.

\remark{The English original of this problem uses the U.S.\ three-year MACRS schedule (33.33\%,
44.45\%, 14.81\%, 7.41\%). With that schedule (book value also zero at $t=4$) the NPV is
\$337{,}231.54: the faster write-off brings tax savings forward, but the decision is unchanged.}

\intuition{The land is not free because the firm already owns it: giving up \$1.5 million today
to get \$1.6 million back in four years costs the project about \$519{,}000 in present value at
13\%. The project still pays for that and for the equipment, so it creates value.}

\refline{Ross, \S10.2 (incremental cash flows: sunk and opportunity costs, NWC), pp.~320--322;
\S10.3--10.4 (pro forma statements, depreciation, salvage), pp.~322--339.}

% =============================================================================
\section{State prices from a traded project --- Amaro de Matos 1.2.6 \#5, p.~28}

\begin{notation}
  \item States $s\in\{1,2\}$ at $t=1$; $\psi=(\psi_1,\psi_2)^\top$: state prices (value today of one unit paid in state $s$ only).
  \item Risk-free asset (asset 0): price 1, payoff $1+r=1.01$ in each state, $r=0.01$.
  \item Project 1: market value $V_1=6.7$, payoffs $(\Gamma_{11},\Gamma_{12})=(10,5)$, investment $I_1=5$.
  \item Project 2: payoffs $(\Gamma_{21},\Gamma_{22})=(-2,4)$, investment $I_2=2$, value $V_2$ to be found.
  \item $\Gamma$: payoff matrix (rows = assets, columns = states); $V$: vector of asset values; $q_s=(1+r)\psi_s$: risk-neutral probabilities.
\end{notation}

\step{1}{The market is complete} The traded assets are the risk-free asset and project 1:
\begin{equation}\label{eq:14gamma}
  \Gamma = \begin{pmatrix} 1.01 & 1.01\\ 10 & 5\end{pmatrix}, \qquad V = \begin{pmatrix} 1\\ 6.7\end{pmatrix}.
\end{equation}
Its determinant is $(1.01)(5)-(1.01)(10)=-5.05\ne0$, so the two payoff vectors are linearly
independent and span both states. The number of independent assets equals the number of states,
so the market is complete (Amaro de Matos \S1.2.1) and the state-price vector, if it exists, is
unique: $\psi=\Gamma^{-1}V$ (\S1.2.2).

\step{2}{No arbitrage gives the pricing equations} By Proposition~1.2.1, absence of arbitrage is
equivalent to a strictly positive $\psi$ with $V=\Gamma\psi$. Row by row:
\begin{align}
  1   &= 1.01\,\psi_1 + 1.01\,\psi_2, \label{eq:14b}\\
  6.7 &= 10\,\psi_1 + 5\,\psi_2. \label{eq:14p}
\end{align}

\step{3}{(a) Solve for $\psi$} Dividing \eqref{eq:14b} by 1.01:
\begin{equation}\label{eq:14sum}
  \psi_1 + \psi_2 = \frac{1}{1.01} = 0.990099.
\end{equation}
So $\psi_2 = 0.990099-\psi_1$. Substituting into \eqref{eq:14p}:
\begin{align}
  6.7 &= 10\psi_1 + 5(0.990099-\psi_1) \notag\\
      &= 5\psi_1 + 4.950495, \notag\\
  \psi_1 &= \frac{6.7-4.950495}{5} = \frac{1.749505}{5} \notag\\
         &= 0.349901. \label{eq:14psi1}
\end{align}
Substituting \eqref{eq:14psi1} into \eqref{eq:14sum}:
\begin{align}
  \psi_2 &= 0.990099 - (0.349901) \notag\\
         &= 0.640198. \label{eq:14psi2}
\end{align}
Both components are strictly positive, so by Proposition~1.2.1 (and exercise \#1 of the same
section) the prices $V$ admit no arbitrage, and $\psi=(0.349901,\ 0.640198)$ is the unique
state-price vector. Check: $10(0.349901)+5(0.640198)=3.49901+3.20099=6.70000$. Equivalently,
the risk-neutral probabilities are $q_1=(1.01)(0.349901)=0.3534$ and $q_2=0.6466$.

\step{4}{(b) Value of project 2} In a complete market every payoff is priced by $\psi$:
\begin{align}
  V_2 &= \Gamma_{21}\psi_1 + \Gamma_{22}\psi_2 \notag\\
      &= (-2)(0.349901) + (4)(0.640198) \notag\\
      &= -0.699802 + 2.560792 \notag\\
      &= 1.860990. \label{eq:14v2}
\end{align}
Justification by replication: a portfolio of $a$ units of project 1 and $b$ units of the bond
pays project 2's payoffs if $10a+1.01b=-2$ and $5a+1.01b=4$. Subtracting, $5a=-6$, so $a=-1.2$,
and then $b=(4-5(-1.2))/1.01=10/1.01=9.900990$. Its cost is
$(-1.2)(6.7)+(9.900990)(1)=-8.04+9.900990=1.860990$, equal to \eqref{eq:14v2}. If project 2's
payoff traded at any other price, buying the cheaper of the two and selling the dearer would be
an arbitrage, which Step~2 rules out.

\step{5}{NPVs} The NPV of a project is the value it creates today, $V_j-I_j$ (\S1.2.3):
\begin{align}
  \mathrm{NPV}_1 &= (6.7) - (5) = 1.7 > 0, \label{eq:14n1}\\
  \mathrm{NPV}_2 &= (1.860990) - (2) = -0.139010 < 0. \label{eq:14n2}
\end{align}

\step{6}{Every agent invests in project 1 and none in project 2} Take any agent with utility
increasing in consumption.
(i) Project 2: whatever future consumption the agent obtains by investing 2 in project 2 can be
obtained by buying the replicating portfolio of Step~4 for 1.860990. The second route leaves
$0.139010$ more for consumption today with the same future consumption, so by monotonicity it
is strictly preferred; no agent invests in project 2.
(ii) Project 1: investing 5 and selling claims to its payoff at the market value 6.7 raises
today's consumption by 1.7 and leaves future consumption unchanged, which every agent strictly
prefers. Hence agents invest in project 1 only, and this does not depend on their preferences.

\intuition{With two states and two independent traded assets, every state-contingent payoff has
one market price. A project is worth doing if it produces its payoff more cheaply than the
market sells it: project 1 does (5 against 6.7), project 2 does not (2 against 1.86).}

\bookcheck{Matches the book's answer (p.~247): $\psi\approx(0.35,\ 0.64)$, $V_2=1.86$,
$\mathrm{NPV}_1=1.7$, $\mathrm{NPV}_2=-0.14$, invest in project 1 only.}

\refline{Amaro de Matos, \S1.2.1 (one-period model, completeness), pp.~18--21; \S1.2.2
(Proposition~1.2.1), pp.~22--23; \S1.2.3 (arbitrage opportunities and investment), pp.~23--24;
\S1.2.4 (risk-neutral measure), p.~24.}

% =============================================================================
\section{Investment independent of preferences under uncertainty --- Amaro de Matos 1.2.6 \#4, p.~28}

\begin{notation}
  \item Agents $i\in\{A,B\}$; dates $t=0,1$; states $s\in\{1,2\}$ at $t=1$ with probabilities $p_1=p$, $p_2=1-p$, $0<p<1$.
  \item $X_i>0$: endowment at $t=0$; $Y_i^s\ge0$: endowment at $t=1$ in state $s$ (it may differ across states).
  \item $\epsilon_i\ge0$: productive investment at $t=0$; $F_s(\epsilon)$: its output at $t=1$ in state $s$ (same technology for both agents).
  \item $x_i$: consumption at $t=0$; $y_i^s$: consumption at $t=1$ in state $s$.
  \item $U_i(x,y)$: utility with the properties of \S1.1 (continuous, twice differentiable, $U_x,U_y>0$, concave);
        $V_i=p\,U_i(x,y^1)+(1-p)\,U_i(x,y^2)$.
  \item $r$: risk-free rate; $\psi_s>0$: state prices; $z_i^s$: holding of the Arrow security paying 1 in state $s$ only (any sign).
  \item $\Pi(\epsilon)\equiv\psi_1F_1(\epsilon)+\psi_2F_2(\epsilon)-\epsilon$: market value created by investing $\epsilon$.
\end{notation}

\noindent\textbf{Claim.} The following conditions are sufficient for both agents to choose the
same investment $\epsilon^\ast$, independent of $U_A,U_B$ and of the endowments:
\begin{itemize}[leftmargin=2.2em]
  \item[(S1)] \emph{Complete, frictionless markets at $t=0$:} besides the risk-free bond there
  is a traded security whose payoff is not proportional to $(1,1)$, so that (with the bond)
  any state-contingent payoff can be bought; no short-sale limits; agents take prices as given.
  Equivalently, the two Arrow securities trade at prices $\psi_1,\psi_2$.
  \item[(S2)] \emph{No arbitrage:} $\psi_1,\psi_2>0$ (Proposition~1.2.1); the bond is priced
  consistently, $\psi_1+\psi_2=1/(1+r)$.
  \item[(S3)] \emph{Technology:} $F_1,F_2$ are differentiable, nondecreasing and concave, at least
  one strictly concave, and there is $\bar\epsilon$ with $\psi_1F_1'(\bar\epsilon)+\psi_2F_2'(\bar\epsilon)<1$
  (the analogue of Assumption~1.1.5: beyond some scale, investing returns less than the market).
\end{itemize}
Then $\epsilon^\ast$ is the unique maximiser of $\Pi$ on $[0,\infty)$, characterised by
\begin{equation}\label{eq:15foc}
  \psi_1F_1'(\epsilon^\ast)+\psi_2F_2'(\epsilon^\ast)=1 \quad\text{if } \psi_1F_1'(0)+\psi_2F_2'(0)>1,
  \qquad \epsilon^\ast=0 \text{ otherwise.}
\end{equation}

\step{1}{Budget constraints} Agent $i$ consumes, invests and buys Arrow securities at $t=0$, and
consumes endowment, output and security payoffs at $t=1$:
\begin{align}
  x_i + \epsilon_i + \psi_1 z_i^1 + \psi_2 z_i^2 &= X_i, \label{eq:15b0}\\
  y_i^s &= Y_i^s + F_s(\epsilon_i) + z_i^s, \qquad s=1,2. \label{eq:15b1}
\end{align}
Lending $b$ at the risk-free rate is the case $z_i^1=z_i^2=(1+r)b$, which by (S2) costs
$(\psi_1+\psi_2)(1+r)b=b$; so the bond is included in \eqref{eq:15b0}--\eqref{eq:15b1}.

\step{2}{Collapse to one wealth constraint} From \eqref{eq:15b1}, $z_i^s = y_i^s - Y_i^s - F_s(\epsilon_i)$.
Substituting into \eqref{eq:15b0}:
\begin{align}
  x_i + \epsilon_i + \psi_1\big(y_i^1 - Y_i^1 - F_1(\epsilon_i)\big) + \psi_2\big(y_i^2 - Y_i^2 - F_2(\epsilon_i)\big) &= X_i. \notag
\end{align}
Moving the terms that do not involve consumption to the right:
\begin{align}
  x_i + \psi_1 y_i^1 + \psi_2 y_i^2 &= X_i + \psi_1Y_i^1 + \psi_2Y_i^2 + \psi_1F_1(\epsilon_i) + \psi_2F_2(\epsilon_i) - \epsilon_i \notag\\
  &= \underbrace{X_i + \psi_1Y_i^1 + \psi_2Y_i^2}_{\text{endowment wealth } W_i^0} + \Pi(\epsilon_i). \label{eq:15w}
\end{align}
Because $z_i^1,z_i^2$ are unrestricted (S1), for a given $\epsilon_i$ every consumption plan
satisfying \eqref{eq:15w} can be reached (choose $z_i^s$ from \eqref{eq:15b1}), and every
reachable plan satisfies \eqref{eq:15w}. So the agent's choice set, given $\epsilon_i$, is the
budget line \eqref{eq:15w} with wealth $W_i^0+\Pi(\epsilon_i)$.

\step{3}{Any optimal investment maximises $\Pi$} Let $\epsilon^\ast$ maximise $\Pi$ (it exists by
Step~4) and suppose an agent's optimum used some $\epsilon'$ with $\Pi(\epsilon')<\Pi(\epsilon^\ast)$,
with consumption $(x',y'^1,y'^2)$ satisfying \eqref{eq:15w} at $\epsilon'$. Define
\begin{equation}\label{eq:15dev}
  x'' = x' + \big(\Pi(\epsilon^\ast)-\Pi(\epsilon')\big) > x'.
\end{equation}
Substituting $(x'',y'^1,y'^2)$ and $\epsilon^\ast$ into the left side of \eqref{eq:15w}:
\begin{align}
  x'' + \psi_1y'^1 + \psi_2y'^2 &= \big(x' + \Pi(\epsilon^\ast)-\Pi(\epsilon')\big) + \psi_1y'^1 + \psi_2y'^2 \notag\\
  &= \big(W_i^0 + \Pi(\epsilon')\big) + \Pi(\epsilon^\ast) - \Pi(\epsilon') \notag\\
  &= W_i^0 + \Pi(\epsilon^\ast), \notag
\end{align}
so the plan is feasible with $\epsilon^\ast$. Since $U_x>0$ (Assumption~1.1.2),
$V_i(x'',y'^1,y'^2)>V_i(x',y'^1,y'^2)$, contradicting optimality. Hence every agent's optimal
investment lies in $\arg\max_{\epsilon\ge0}\Pi(\epsilon)$. The argument used only that utility
increases in $x$: not the form of $U_i$, not $(X_i,Y_i^1,Y_i^2)$, and not even $p$.

\step{4}{The maximiser is unique and solves \eqref{eq:15foc}} Differentiating $\Pi$:
\begin{align}
  \Pi'(\epsilon) &= \psi_1F_1'(\epsilon)+\psi_2F_2'(\epsilon)-1, \label{eq:15d1}\\
  \Pi''(\epsilon) &= \psi_1F_1''(\epsilon)+\psi_2F_2''(\epsilon) < 0, \label{eq:15d2}
\end{align}
where \eqref{eq:15d2} uses $\psi_s>0$ (S2) and concavity with one $F_s$ strictly concave (S3).
So $\Pi$ is strictly concave and has at most one maximiser. By (S3), $\Pi'(\bar\epsilon)<0$, and
$\Pi'$ is decreasing, so $\Pi$ falls beyond $\bar\epsilon$: a maximiser exists on the compact set
$[0,\bar\epsilon]$. The Kuhn--Tucker conditions for $\max_{\epsilon\ge0}\Pi$ are $\Pi'(\epsilon^\ast)\le0$,
$\epsilon^\ast\ge0$, $\epsilon^\ast\Pi'(\epsilon^\ast)=0$; with \eqref{eq:15d1} they give
\eqref{eq:15foc}. Both agents therefore choose $\epsilon_A=\epsilon_B=\epsilon^\ast$. $\square$

\step{5}{Link with the agents' first-order conditions (the book's route)} Substituting
\eqref{eq:15b0}--\eqref{eq:15b1} into $V_i$ and differentiating with respect to $z_i^s$ and
$\epsilon_i$ at an interior optimum:
\begin{align}
  \frac{\partial V_i}{\partial z_i^s} &= -\psi_s\,\bar U_x + p_s\,U_y^s = 0
     \ \Longrightarrow\ \psi_s = \frac{p_s\,U_y^s}{\bar U_x}, \label{eq:15fz}\\
  \frac{\partial V_i}{\partial \epsilon_i} &= -\bar U_x + p\,U_y^1F_1'(\epsilon_i) + (1-p)\,U_y^2F_2'(\epsilon_i) = 0, \label{eq:15fe}
\end{align}
where $\bar U_x\equiv p\,\partial_xU_i(x_i,y_i^1)+(1-p)\,\partial_xU_i(x_i,y_i^2)>0$ and $U_y^s\equiv\partial_yU_i(x_i,y_i^s)>0$.
Dividing \eqref{eq:15fe} by $\bar U_x$ and substituting \eqref{eq:15fz}:
\begin{align}
  \Big(\frac{p\,U_y^1}{\bar U_x}\Big)F_1'(\epsilon_i) + \Big(\frac{(1-p)\,U_y^2}{\bar U_x}\Big)F_2'(\epsilon_i) &= 1, \notag\\
  (\psi_1)F_1'(\epsilon_i) + (\psi_2)F_2'(\epsilon_i) &= 1, \label{eq:15same}
\end{align}
which is \eqref{eq:15foc}. Each agent's marginal rates of substitution are individual, but in a
complete market each agent trades until they equal the \emph{common} prices $\psi_s$, and only
those prices enter \eqref{eq:15same}. With $q_s=(1+r)\psi_s$ (so $q_1+q_2=1$ by (S2)),
\eqref{eq:15same} reads $q_1F_1'(\epsilon^\ast)+q_2F_2'(\epsilon^\ast)=1+r$: the expected marginal
return under the risk-neutral measure equals the gross risk-free rate, the uncertainty version
of $\phi(\epsilon^\ast)=r$ in \S1.1.5.

\step{6}{Why completeness is needed, and a second sufficient condition} With only the bond,
$z_i^1=z_i^2$ is forced and \eqref{eq:15fz} is replaced by the single condition
$\bar U_x=(1+r)\big(pU_y^1+(1-p)U_y^2\big)$. Dividing \eqref{eq:15fe} by it gives
\begin{equation}\label{eq:15inc}
  m_iF_1'(\epsilon_i)+(1-m_i)F_2'(\epsilon_i)=1+r, \qquad
  m_i\equiv\frac{p\,U_y^1}{p\,U_y^1+(1-p)\,U_y^2},
\end{equation}
and $m_i$ depends on agent $i$'s utility and, through $y_i^s$, on endowments: investment is no
longer separated from preferences. The exception is when the marginal return carries no
uncertainty, $F_1'(\epsilon)=F_2'(\epsilon)\equiv F'(\epsilon)$ for all $\epsilon$ (the book's
``trivial uncertainty''). Substituting into \eqref{eq:15inc}:
\begin{align}
  m_i\,(F'(\epsilon_i)) + (1-m_i)\,(F'(\epsilon_i)) &= 1+r, \notag\\
  F'(\epsilon_i) &= 1+r, \notag
\end{align}
the same for every agent. So either (S1)--(S3), or a bond market plus nonstochastic marginal
returns (with concave $F$), is sufficient.

\intuition{Fisher separation survives uncertainty when markets let every agent reshape the
project's risky output into whatever consumption pattern they like at common prices. Then the
best project is the one with the highest market value, and all agents, however patient or risk
averse, agree on it. Without complete markets each agent values the risky output with private
marginal utilities, and investment depends on who is investing.}

\bookcheck{Matches the book's answer (pp.~245--247): a unique state-price vector (two states and
two independent assets) gives an investment rule, its (C.3), $q(\theta_1)F_\epsilon(\epsilon,\theta_1)+q(\theta_2)F_\epsilon(\epsilon,\theta_2)=1+r$,
that involves neither $U$, $x$, $y$ nor the endowments. The book identifies its weights with
$p_sU_y^s$; they equal \eqref{eq:15same}'s $q_s=(1+r)\psi_s$ once normalised to sum to one. The
solution above adds the concavity and scale conditions (S3) that make $\epsilon^\ast$ exist and be
unique, and proves the result without differentiating $U$ (Step~3).}

\refline{Amaro de Matos, \S1.1.1 (Assumptions 1.1.1--1.1.3), pp.~5--6; \S1.1.5 (Proposition~1.1.1,
Fisher separation), pp.~11--13; \S1.2.1--1.2.3 (state prices, completeness, Proposition~1.2.1,
arbitrage and investment), pp.~18--24.}

\end{document}
